\section{Условие}
{\bfseries Цель работы:} приобретение практических навыков в управлении потоками в ОС и в обеспечении синхронизации между потоками.

{\bfseries Задание:} составить программу на языке Си, обрабатывающую данные в многопоточном режиме. При обработке использовать стандартные средства создания потоков операционной системы (Windows/Unix). Ограничение максимального количества потоков, работающих в один момент времени, должно быть задано ключом запуска программы. Так же необходимо уметь продемонстрировать количество потоков, используемое программой, с помощью стандартных средств операционной системы. В отчёте привести исследование зависимости ускорения и эффективности алгоритма от входных данных и количества потоков. Получившиеся результаты необходимо объяснить.

{\bfseries Вариант:} 20. Дан массив координат $(x, y, z)$. Необходимо найти три точки, которые образуют треугольник максимальной площади.
